\documentclass[10pt,a4paper]{article}
\usepackage[a4paper,margin=14mm,top=12mm,bottom=14mm]{geometry}
\usepackage{fontspec}
\setmainfont{Liberation Sans}
\usepackage{polyglossia}
\setdefaultlanguage{french}
\usepackage{amsmath}
\usepackage{xcolor}
\usepackage{graphicx}
\usepackage[most]{tcolorbox}
\usepackage{tabularx,booktabs,array,multirow}
\usepackage{enumitem}
\usepackage{fancyhdr}
\usepackage{lastpage}
\usepackage[hidelinks]{hyperref}
\definecolor{uimm}{HTML}{005677}
\definecolor{cyan}{HTML}{13A2DD}
\definecolor{fond}{HTML}{E6F2FA}
\definecolor{ok}{HTML}{2E7D32}
\definecolor{ko}{HTML}{B3261E}
\definecolor{violet}{HTML}{7B1FA2}
\setlength{\parindent}{0pt}
\setlength{\parskip}{3pt}
\setlist{nosep}
\setlist[itemize]{leftmargin=13pt,label=\textcolor{cyan}{\textbullet}}
\newcommand{\titre}[1]{\vspace{5pt}{\color{uimm}\large\bfseries #1}\par\vspace{-3pt}{\color{cyan}\rule{\linewidth}{1.2pt}}\par\vspace{0pt}}
\tcbset{boite/.style={colback=fond,colframe=uimm,boxrule=0.6pt,arc=2pt,left=6pt,right=6pt,top=4pt,bottom=4pt},
        defi/.style={colback=violet!6,colframe=violet,boxrule=0.8pt,arc=2pt,left=6pt,right=6pt,top=4pt,bottom=4pt,
                     title={\bfseries Défi},coltitle=white,colbacktitle=violet,fonttitle=\small}}
\newcommand{\bandeau}[3]{%
\begin{tcolorbox}[colback=uimm,colframe=uimm,arc=2pt,left=4pt,right=8pt,top=4pt,bottom=4pt]
\begin{minipage}[c]{18mm}\setlength{\fboxsep}{1.2pt}\fcolorbox{white}{white}{\includegraphics[height=16.5mm]{logo_uimm.jpg}}\end{minipage}\hspace{4mm}%
\begin{minipage}[c]{\dimexpr\linewidth-22mm\relax}\color{white}
{\small #1}\\[1pt]{\LARGE\bfseries #2}\\[2pt]{\large #3}\hfill{\small CQPM Technicien qualité}
\end{minipage}
\end{tcolorbox}}
\newcommand{\reponse}[1][1]{\par\vspace{2pt}{\color{gray!60}\foreach \i in {1,...,#1}{\rule{\linewidth}{0.3pt}\par\vspace{9pt}}}}
\usepackage{pgffor}
\pagestyle{fancy}\fancyhf{}
\renewcommand{\headrulewidth}{0pt}
\newcommand{\pied}[1]{\fancyfoot[L]{\footnotesize\color{uimm}#1}\fancyfoot[R]{\footnotesize\color{uimm}Page \thepage{} sur \pageref{LastPage}}}

\usepackage{tikz}
\usepackage{longtable}
\renewcommand{\bandeau}[3]{%
\begin{tcolorbox}[colback=uimm,colframe=uimm,arc=2pt,left=4pt,right=8pt,top=4pt,bottom=4pt]
\begin{minipage}[c]{18mm}\setlength{\fboxsep}{1.2pt}\fcolorbox{white}{white}{\includegraphics[height=16.5mm]{logo_uimm.jpg}}\end{minipage}\hspace{4mm}%
\begin{minipage}[c]{\dimexpr\linewidth-22mm\relax}\color{white}
{\small #1}\\[1pt]{\LARGE\bfseries #2}\\[2pt]{\large #3}
\end{minipage}
\end{tcolorbox}}
% en-tête d'une fiche : code, titre, question, bloc, durée, organisation, niveau
\newcommand{\entete}[7]{%
\begin{tcolorbox}[colback=uimm,colframe=uimm,arc=2pt,left=4pt,right=6pt,top=4pt,bottom=4pt]
\begin{minipage}[c]{18mm}\setlength{\fboxsep}{1.2pt}\fcolorbox{white}{white}{\includegraphics[height=16.5mm]{logo_uimm.jpg}}\end{minipage}\hspace{3mm}%
\begin{minipage}[c]{\dimexpr\linewidth-58mm\relax}\color{white}\raggedright
{\small Base SPC · #4}\\[1pt]{\Large\bfseries #1 · #2}\\[2pt]{\itshape #3}
\end{minipage}\hfill
\begin{minipage}[c]{33mm}\color{white}\small\raggedleft
Durée : \textbf{#5}\\ \textbf{#6}\\ Niveau #7 sur 3
\end{minipage}
\end{tcolorbox}}
\newcommand{\tagE}{{\setlength{\fboxsep}{1.5pt}\colorbox{cyan}{\color{white}\scriptsize\bfseries Ellistat}}}
\newcommand{\tagW}{{\setlength{\fboxsep}{1.5pt}\colorbox{uimm}{\color{white}\scriptsize\bfseries Page web}}}
\tcbset{pred/.style={colback=orange!6,colframe=orange!80!black,boxrule=0.6pt,arc=2pt,left=6pt,right=6pt,top=3pt,bottom=1pt,
                    title={\bfseries Je prédis, avant de commencer},coltitle=white,colbacktitle=orange!80!black,fonttitle=\small},
        retenir/.style={colback=fond,colframe=uimm,boxrule=0.6pt,arc=2pt,left=6pt,right=6pt,top=3pt,bottom=3pt,
                    title={\bfseries À retenir},coltitle=white,colbacktitle=uimm,fonttitle=\small},
        plus/.style={colback=violet!6,colframe=violet,boxrule=0.6pt,arc=2pt,left=6pt,right=6pt,top=3pt,bottom=3pt,
                    title={\bfseries Pour aller plus loin},coltitle=white,colbacktitle=violet,fonttitle=\small},
        act/.style={colback=uimm!10,colframe=uimm,boxrule=0.5pt,arc=2pt,left=5pt,right=5pt,top=2pt,bottom=2pt,before skip=7pt,after skip=3pt},
        verif/.style={colback=orange!6,colframe=orange!80!black,boxrule=0.6pt,arc=2pt,left=6pt,right=6pt,top=3pt,bottom=3pt,
                    title={\bfseries À vérifier sur poste avant la première séance},coltitle=white,colbacktitle=orange!80!black,fonttitle=\small}}
\newcommand{\fichier}[1]{\texttt{\small #1}}
\newcommand{\menu}[1]{\textsf{\textbf{#1}}}
\newcommand{\champ}[1]{\rule{#1}{0.3pt}}

\pied{Base SPC · fiche M14 · Trois étalons pour la butée}
\begin{document}
\thispagestyle{fancy}
\entete{M14}{Trois étalons pour la butée}{Le micromètre est-il aussi juste aux deux limites de la tolérance qu'en son milieu ?}{Le moyen de mesure}{60 min}{en binôme}{3}
\begin{tcolorbox}[boite]\textbf{Le but.} Composer trois étalons par adhérence (11,850 ; 11,925 ; 12,000 mm), mesurer chacun et vérifier la justesse là où l'on décide : aux deux limites et au milieu de la tolérance.\par\textbf{Il vous faut.} Un coffret de cales (cales de 0,5 ; 1,005 ; 1,35 ; 1,42 ; 2 ; 9,5 et 10 mm) · Un micromètre 0-25 mm · Chiffon non pelucheux · Ellistat\end{tcolorbox}
\begin{tcolorbox}[pred]Votre micromètre aura-t-il le même écart de justesse aux deux limites de la tolérance et en son milieu ? Écrivez l'écart que vous attendez, en µm.\reponse\end{tcolorbox}
\titre{À faire}
\begin{enumerate}[label=\textcolor{cyan}{\textbf{\arabic*.}},leftmargin=15pt,itemsep=1.5pt]
\item \textbf{Composer} les trois étalons du tableau avec les cales du coffret. Méthode : on règle d'abord le \textbf{dernier chiffre}. Pour 11,925, la seule cale qui donne le 0,005 est la 1,005 : il reste 10,920 à obtenir avec deux cales. Vérifier chaque somme.
\item Nettoyer les faces. Coller les cales deux à deux, en croix, en pressant et en glissant, comme en M10. Ne pas tenir les piles dans la main.
\item Mesurer \textbf{10 fois} chaque pile au cliquet, en l'enlevant à chaque fois, et en \textbf{alternant} : 11,850 ; 11,925 ; 12,000 ; puis on recommence. La température agit ainsi de la même façon sur les trois.
\item \tagE\ \menu{Gage R\&R}, onglet \menu{Justesse} : trois études (« Etalon\_11\_850 », « Etalon\_11\_925 », « Etalon\_12\_000 »), comme en M3. Nombre de mesures 10 ; Valeur vraie = valeur de l'étalon ; LSI 11,85 et LSS 12,00 ; k 3. Relever Cg, Cgk et le Test T.
\item À la main, pour chaque étalon : le biais en µm et le rapport $|\text{biais}| / s$. Pour 10 mesures, le biais est négligeable si ce rapport est \textbf{au plus 0,715} (livret Bosch, annexe C).
\item Séparer les cales en les faisant glisser, sans tirer. Les nettoyer et les ranger.
\end{enumerate}
\vspace{3pt}{\renewcommand{\arraystretch}{1.45}\begin{tabularx}{\linewidth}{|c|X|c|}\hline
\textbf{Étalon visé} & \textbf{Cales choisies (mm)} & \textbf{Somme (mm)}\\\hline
11,850 mm & & \\\hline
11,925 mm & & \\\hline
12,000 mm & & \\\hline
\end{tabularx}}\par\vspace{5pt}
{\renewcommand{\arraystretch}{1.15}\small\begin{tabularx}{\linewidth}{|l|*{3}{>{\centering\arraybackslash}X|}}\hline
\textbf{Mesure} & \textbf{11,850} & \textbf{11,925} & \textbf{12,000}\\\hline
1 & & & \\\hline
2 & & & \\\hline
3 & & & \\\hline
4 & & & \\\hline
5 & & & \\\hline
6 & & & \\\hline
7 & & & \\\hline
8 & & & \\\hline
9 & & & \\\hline
10 & & & \\\hline
\textbf{Moyenne (mm)} & & & \\\hline
\textbf{Biais = moyenne $-$ étalon (µm)} & & & \\\hline
\textbf{$|\text{biais}| / s$ (à la main)} & & & \\\hline
\textbf{Cg (Ellistat)} & & & \\\hline
\textbf{Cgk (Ellistat)} & & & \\\hline
\end{tabularx}}\par
\newpage
\titre{J'observe, je décide}
\begin{enumerate}[label=\textcolor{cyan}{\textbf{Q\arabic*.}},leftmargin=22pt]
\item Écart entre le plus grand et le plus petit biais : \champ{16mm} µm. Le biais est-il le même aux trois étalons ?\reponse[1]
\item Pour quel(s) étalon(s) le rapport $|\text{biais}| / s$ dépasse-t-il 0,715 ? Faut-il en conclure que le micromètre est déréglé à cet endroit ?\reponse[2]
\item Le micromètre est-il capable aux deux limites de la tolérance ? Appuyez-vous sur Cgk.\reponse[2]
\item Chaque cale d'une pile a sa propre incertitude, de l'ordre de quelques dixièmes de micromètre. Cela change-t-il vos conclusions, avec un micromètre qui lit 1 µm ?\reponse[2]
\end{enumerate}
\begin{tcolorbox}[retenir]On qualifie un instrument là où il sert : aux deux limites et au milieu de la tolérance. Un seuil statistique n'est pas une frontière : un écart de 0,2 µm entre deux étalons ne change pas l'instrument.\end{tcolorbox}
\begin{tcolorbox}[plus]Dans Ellistat, ouvrez \menu{Linéarité} avec les trois étalons (3 produits, 10 valeurs). Les trois points couvrent 0,15 mm seulement : lisez la pente avec prudence. Ou complétez l'étalon de 11,925 à 25 mesures, comme le demande le livret Bosch (§ 4.1).\end{tcolorbox}
\vfill{\footnotesize\color{uimm}\textbf{Ressources de la base :} Données de secours · L'instrument est-il capable ?.}
\end{document}